\sectionkicker{Theme synthesis}
\subsection*{横断比較 — 6月のモデル競争は「何を出したか」だけでは測れない}
\addcontentsline{toc}{subsection}{横断比較 — 6月のモデル競争は「何を出したか」だけでは測れない}
6月の主要releaseを並べると、model capabilityの主張そのものより、open-weightかmanaged APIか、multimodal/computer-use/codingへどう接続したか、preview・GA・deprecationを含む提供状態が採用判断を左右していたことが見える。
\medskip
\begin{center}
\small
\renewcommand{\arraystretch}{1.16}
\begin{tabularx}{\linewidth}{@{}p{0.20\linewidth}X@{}}
\toprule
観察軸 & 6月の一次資料・論文から読めること \\
\midrule
配布・提供形態 & MiniMax M3はopen weightsを伴うreleaseとして提示された一方、Anthropic、Alibaba Model Studio、Gemini APIはmanaged product/API lifecycleの更新として確認される。Kimi K2.7 Codeはcode modelのrelease/open-sourceとKimi Code側のworkflow更新が接続した。 \cite{src-d8e13386974c,src-f6aa679babd7,src-eac8dfbeb78c,src-a24fd97eb3d0,src-9191d43918fc} \\
agent / computer-use接続 & MiniMax M3はdesktop operationをvendor claimとして掲げ、Gemini APIではComputer Use、Kimi CodeではGoal modeと/swarmが同月に更新された。モデル単体の公開と、実行環境へ接続する機能が同じ競争面へ寄っている。 \cite{src-d8e13386974c,src-a24fd97eb3d0,src-9191d43918fc} \\
availability / lifecycle & Geminiではpreview、GA、deprecation、shutdownが同じ月次changelogに並び、Alibabaでもsnapshotやregion別availabilityが更新された。GPT-5.6は6月26日時点ではlimited previewであり、翌月の一般的な展開と同一視しないことが重要である。 \cite{src-a24fd97eb3d0,src-eac8dfbeb78c,src-30f48121decc} \\
modalitiesと用途 & MiniMax M3のimage/video input、Alibabaのvisual/video系更新、Google DeepMindのDiffusionGemma等、Kimiのcoding workflowは同じ『frontier』でも用途面が異なる。単一benchmark順位だけでは6月の差分を表しにくい。 \cite{src-d8e13386974c,src-eac8dfbeb78c,src-dd5d9b8662f3,src-9191d43918fc} \\
\bottomrule
\end{tabularx}
\renewcommand{\arraystretch}{1.0}
\end{center}
\normalsize
\begin{claimboundary}[この比較の境界]
この表は本号で既に検証・参照している一次資料と論文を横断比較の形へ再配置したものである。vendor / project / author が報告した性能値や優位性は独立再現済みの一般則へ変換せず、本文とSource Notesの帰属境界をそのまま維持する。
\end{claimboundary}
